% ============================================================
\chapter{Modul 13: Matplotlib -- Daten visualisieren}
% ============================================================

\section{Lektion 36: Diagramme erstellen}

\begin{lstlisting}
import numpy as np
import matplotlib.pyplot as plt

# Simulierte Sensordaten
zeit       = np.linspace(0, 10, 100)
abstand    = 1.5 + 0.5*np.sin(zeit) + np.random.normal(0, 0.05, 100)
akku       = 100 - zeit * 3.5

fig, (ax1, ax2) = plt.subplots(2, 1, figsize=(10, 6))
fig.suptitle("MobileRobot-1: Testfahrt-Auswertung", fontsize=14)

# Abstandsverlauf
ax1.plot(zeit, abstand, 'b-', linewidth=1.5, label="Abstand roh")
ax1.axhline(y=0.25, color='r', linestyle='--', label="Gefahrenschwelle")
ax1.fill_between(zeit, abstand, 0.25,
                 where=(abstand < 0.25), alpha=0.3, color='red',
                 label="Gefahrenbereich")
ax1.set_ylabel("Abstand [m]")
ax1.legend()
ax1.grid(True, alpha=0.3)

# Akku-Verlauf
ax2.plot(zeit, akku, 'g-', linewidth=2)
ax2.axhline(y=20, color='orange', linestyle='--', label="Warnschwelle")
ax2.set_xlabel("Zeit [s]")
ax2.set_ylabel("Akku [%]")
ax2.legend()
ax2.grid(True, alpha=0.3)

plt.tight_layout()
plt.savefig("testfahrt.pdf", dpi=150)
plt.show()
\end{lstlisting}

\section{Lektion 37: Roboter-Trajektorie und Polar-Plot}

\begin{lstlisting}
# Trajektorie des Roboters
theta_verlauf = np.linspace(0, 2*np.pi, 200)
x_bahn = np.cumsum(np.cos(theta_verlauf)) * 0.05
y_bahn = np.cumsum(np.sin(theta_verlauf)) * 0.05

fig, (ax1, ax2) = plt.subplots(1, 2, figsize=(12, 5))

# Fahrspur
ax1.plot(x_bahn, y_bahn, 'b-', linewidth=1)
ax1.plot(x_bahn[0], y_bahn[0], 'gs', ms=10, label="Start")
ax1.plot(x_bahn[-1], y_bahn[-1], 'r*', ms=10, label="Ende")
ax1.set_aspect('equal')
ax1.legend()
ax1.set_title("Roboter-Trajektorie")
ax1.grid(True)

# Laserscanner: Polar-Plot
winkel = np.linspace(0, 2*np.pi, 360)
distanzen = 1.5 + 0.8*np.sin(3*winkel) + np.random.uniform(0, 0.1, 360)

ax2 = plt.subplot(1, 2, 2, projection='polar')
ax2.plot(winkel, distanzen, 'r-', linewidth=0.8)
ax2.fill(winkel, distanzen, alpha=0.2, color='red')
ax2.set_title("Laserscanner (360 Grad)")

plt.tight_layout()
plt.savefig("roboter_visualisierung.pdf", dpi=150)
plt.show()
\end{lstlisting}

\begin{praxis}
Visualisierung ist in der Robotik kein Luxus -- es ist Debugging. Ein Polar-Plot des Laserscanners zeigt sofort, wo Hindernisse sind. Eine Trajektorie zeigt, ob der Roboter wirklich geradeaus fährt oder driftet.
\end{praxis}

\begin{uebung}[title=Projektaufgabe Modul 13: Vollständiges Dashboard]
Erstelle ein 4-Panel-Dashboard für eine simulierte Testfahrt: Panel 1 -- Trajektorie (x/y), Panel 2 -- Abstandsverlauf über Zeit, Panel 3 -- Akkuverlauf, Panel 4 -- Geschwindigkeitsverteilung (Histogramm). Alle Panels mit Beschriftung, Legende und Grid.
\end{uebung}
